\writeSectionFrame{\presentationTitle}{Das Problem}
\frameWithImageOnly{\BILDERPFAD/problem1.jpg}
\note{
 	Wenn man sich eine Kaffeemaschine vorstellt, die verschiedene Kaffeesorten zubereiten kann, fällt auf, dass es viele Ähnlichkeiten beim Zubereitungsalgorithmus gibt.
}

\begin{frame}{Kaffeezubereitungsvorgang}
	\begin{itemize}
 		\item Wasser kochen
 		\item Kaffeebohnen mahlen
 		\item Brühen
 		\item In Tasse gießen
 	\end{itemize}
\end{frame}

\begin{frame}{Normaler Kaffee}
	\begin{itemize}
 		\item Wasser kochen
 		\item Kaffeebohnen mahlen
 		\item \textbf{\color{orange}Brühen\color{black}}
 		\item In Tasse gießen
 	\end{itemize}
\end{frame}

\begin{frame}{Arabischer Kaffee}
	\begin{itemize}
 		\item Wasser kochen
 		\item Kaffeebohnen mahlen
 		\item \textbf{\color{orange}Brühen\color{black}}
 		\item In Tasse gießen
 	\end{itemize}
\end{frame}

\begin{frame}{Cappuccino}
	\begin{itemize}
 		\item Wasser kochen
 		\item Kaffeebohnen mahlen
 		\item \textbf{\color{orange}Brühen\color{black}}
        \item \textbf{\color{orange}Milch hinzufügen\color{black}}
 		\item In Tasse gießen
 	\end{itemize}
\end{frame}
\note{
	Die Rezepte sind nahezu identisch und wir müssten viel doppelten Code
	schreiben. Nur der erste markierte Schritt unterscheidet sich. Es wäre doch viel
	einfacher, wenn wir den Algorithmus allgemein gültig definieren und dann den
	ersten Schritt entsprechend austauschen könnten. Wenn wir also eine anpassbare Schablone hätten für unseren Algorithmus.
}

